% ============================================================
\chapter{D\'erivation et Int\'egration Num\'erique}
\label{ch:integration}
% ============================================================

\section{Exemple motivant}
Calculer $\int_0^1 e^{-x^2}\,dx$ : pas de primitive en forme close. On utilise des formules de quadrature. De m\^eme, $f'(x)$ peut \^etre approch\'ee par des diff\'erences finies lorsque $f$ n'est connue que num\'eriquement.

\section{D\'erivation num\'erique}\index{derivation numerique@d\'erivation num\'erique}

\begin{definition}[Diff\'erences finies]
\begin{align}
f'(x) &\approx \frac{f(x+h)-f(x)}{h} && \text{(progressive, } \BigO(h)\text{)} \\
f'(x) &\approx \frac{f(x+h)-f(x-h)}{2h} && \text{(centr\'ee, } \BigO(h^2)\text{)} \\
f''(x) &\approx \frac{f(x+h)-2f(x)+f(x-h)}{h^2} && (\BigO(h^2))
\end{align}
\end{definition}

\begin{theoreme}[Erreur de la diff\'erence centr\'ee]
Si $f\in C^3$ : $\frac{f(x+h)-f(x-h)}{2h}=f'(x)+\frac{h^2}{6}f'''(\xi)$.
\end{theoreme}

\begin{attention}
Diminuer $h$ am\'eliore l'erreur de troncature mais augmente l'erreur d'arrondi. Le $h$ optimal est $h^*\sim\eps_{\mathrm{mach}}^{1/3}\approx 6\times 10^{-6}$ en double pr\'ecision.
\end{attention}

\section{Formules de Newton--Cotes}\index{Newton--Cotes}

Principe : remplacer $f$ par son polyn\^ome d'interpolation sur $[a,b]$ et int\'egrer.

\subsection{R\`egle du trap\`eze}\index{trapeze@trap\`eze}

\begin{definition}
\[
\int_a^b f(x)\,dx \approx \frac{b-a}{2}[f(a)+f(b)].
\]
Erreur : $-\frac{(b-a)^3}{12}f''(\xi)$.
\end{definition}

\subsection{R\`egle de Simpson}\index{Simpson}

\begin{definition}
\[
\int_a^b f(x)\,dx \approx \frac{b-a}{6}\left[f(a)+4f\!\left(\frac{a+b}{2}\right)+f(b)\right].
\]
Erreur : $-\frac{(b-a)^5}{2880}f^{(4)}(\xi)$.
\end{definition}

\begin{theoreme}[Formule de Simpson --- degr\'e d'exactitude]
Simpson est exacte pour les polyn\^omes de degr\'e $\leq 3$ (un degr\'e de mieux que pr\'evu).
\end{theoreme}

\section{Formules composites}

\begin{definition}[Trap\`eze composite]
$n$ sous-intervalles, $h=(b-a)/n$ :
\[
T_n(f) = \frac{h}{2}\left[f(a)+2\sum_{i=1}^{n-1}f(x_i)+f(b)\right].
\]
Erreur : $-\frac{(b-a)h^2}{12}f''(\xi)=\BigO(h^2)$.
\end{definition}

\begin{definition}[Simpson composite]
$n$ pairs de sous-intervalles, $h=(b-a)/n$ :
\[
S_n(f) = \frac{h}{3}\left[f(a)+4\sum_{\text{imp}}f(x_i)+2\sum_{\text{pair}}f(x_i)+f(b)\right].
\]
Erreur : $\BigO(h^4)$.
\end{definition}

\section{Extrapolation de Richardson}\index{Richardson}

\begin{theoreme}
Si $T(h)=I+c_1h^2+c_2h^4+\cdots$, alors $\frac{4T(h/2)-T(h)}{3}=I+\BigO(h^4)$. C'est la formule de Simpson !
\end{theoreme}

\section{Exemple num\'erique}

$I=\int_0^1 e^{-x^2}\,dx\approx 0{,}746824$.

\begin{center}
\begin{tabular}{cccc}
\toprule
$n$ & Trap\`eze $T_n$ & Simpson $S_n$ & Erreur Simpson \\
\midrule
2 & 0.731531 & 0.747180 & $3{,}6\times10^{-4}$ \\
4 & 0.742984 & 0.746831 & $7{,}0\times10^{-6}$ \\
8 & 0.745866 & 0.746824 & $1{,}1\times10^{-7}$ \\
16 & 0.746584 & 0.746824 & $1{,}7\times10^{-9}$ \\
\bottomrule
\end{tabular}
\end{center}

\section{Impl\'ementation}

\begin{codeblock}[title=\textbf{Python}]
\begin{minted}{python}
import numpy as np

def trapezoidal(f, a, b, n):
    x = np.linspace(a, b, n+1)
    h = (b - a) / n
    return h * (f(a)/2 + np.sum(f(x[1:-1])) + f(b)/2)

def simpson(f, a, b, n):
    assert n % 2 == 0
    x = np.linspace(a, b, n+1)
    h = (b - a) / n
    return h/3 * (f(a) + 4*np.sum(f(x[1::2])) + 2*np.sum(f(x[2:-1:2])) + f(b))

f = lambda x: np.exp(-x**2)
exact = 0.7468241328124271  # scipy.integrate.quad

for n in [2, 4, 8, 16, 32]:
    T = trapezoidal(f, 0, 1, n)
    S = simpson(f, 0, 1, n)
    print(f"n={n:3d}  Trap={T:.10f}  Simp={S:.10f}  "
          f"Err_T={abs(T-exact):.2e}  Err_S={abs(S-exact):.2e}")
\end{minted}
\end{codeblock}

\begin{codeblock}[title=\textbf{Julia}]
\begin{minted}{julia}
function trapezoidal(f, a, b, n)
    h = (b - a) / n
    x = range(a, b, length=n+1)
    return h * (f(a)/2 + sum(f.(x[2:end-1])) + f(b)/2)
end

function simpson_rule(f, a, b, n)
    h = (b - a) / n
    x = range(a, b, length=n+1)
    return h/3 * (f(a) + 4sum(f.(x[2:2:end-1])) + 2sum(f.(x[3:2:end-2])) + f(b))
end

f(x) = exp(-x^2)
for n in [2, 4, 8, 16, 32]
    println("n=$n  Trap=$(trapezoidal(f,0,1,n))  Simp=$(simpson_rule(f,0,1,n))")
end
\end{minted}
\end{codeblock}

\section{Exercices}

\begin{exercice}[$\star$]
Calculer $\int_0^\pi\sin x\,dx$ par trap\`eze et Simpson composites ($n=4$). Comparer \`a la valeur exacte.
\end{exercice}

\begin{exercice}[$\star\star$]
D\'emontrer que Simpson est exacte pour les polyn\^omes de degr\'e 3.
\end{exercice}

\begin{exercice}[$\star\star$]
V\'erifier num\'eriquement l'ordre $\BigO(h^2)$ du trap\`eze et $\BigO(h^4)$ de Simpson par un tableau d'erreurs.
\end{exercice}

\begin{exercice}[$\star\star\star$ -- Projet]
Impl\'ementer l'int\'egration adaptative (Simpson adaptatif) et comparer avec \texttt{scipy.integrate.quad}.
\end{exercice}

\begin{formules}
\begin{align*}
T_n &= \frac{h}{2}[f(a)+2\textstyle\sum f(x_i)+f(b)], \; E=\BigO(h^2) \\
S_n &= \frac{h}{3}[f(a)+4\textstyle\sum_{\text{imp}}+2\textstyle\sum_{\text{pair}}+f(b)], \; E=\BigO(h^4) \\
f'(x) &\approx \frac{f(x+h)-f(x-h)}{2h}, \; E=\BigO(h^2)
\end{align*}
\end{formules}
